\subsection{Finite local preparation of the scalar intervention}
\label{sec:source-scalar-local-preparation}

The coherent intervention can be prepared by a finite local force while the
original field action continues to evolve. This removes the instantaneous
displacement from the conditional instrument construction. The original
vacuum, quantization and admissibility of the force control remain hypotheses;
neither the vacuum nor quantum outcome statistics are selected by this result.

Use the same convention $\hbar=1$. Let $\mathsf A\ge I$ be self-adjoint in the positive mass inner product
$\langle\cdot,\cdot\rangle_M$, and let $v$ be the compact velocity profile.
Starting with the original vacuum at time $-\delta/2$, apply
\[
 H(t)=H_0-\frac1\delta\sum_{i\in\operatorname{supp}v}m_i v_iq_i,
 \qquad -\delta/2\le t\le\delta/2.
\]
All force ports act in parallel. Their shared duration is $\delta$, and their
strength increases as $1/\delta$. Thereafter evolve with $H_0$ and apply the
declared field--pointer instrument.

\begin{proposition}[Preparation with an active free field]
Up to a scalar phase, the pulse prepares at its endpoint the free evolution
for time $\delta/2$ of the coherent momentum displacement with velocity
\[
 v_\delta=\operatorname{sinc}(\delta\sqrt{\mathsf A}/2)v.
\]
Write $\rho_v(t)$ for free evolution from the ideal displaced vacuum at
the virtual reference time zero, and $\rho_\delta(t)$ for the pulse-prepared
state followed by free evolution. For every $t\ge\delta/2$,
\[
 D(\rho_\delta(t),\rho_v(t))
 \le \frac{\delta^2}{24\sqrt2}\|\mathsf A v\|_M,
 \qquad D(\rho,\sigma)=\tfrac12\|\rho-\sigma\|_1.
\]
If the initial state has trace distance at most $\epsilon_{\rm vac}$ from
that vacuum, and an additional real acceleration force $e(t)$ in the same
controlled ports satisfies $\int\|e(t)\|_M\,dt\le\epsilon_F$, then
\[
 D(\widetilde\rho_\delta(t),\rho_v(t))
 \le \epsilon_{\rm vac}+\frac{\epsilon_F}{\sqrt2}
       +\frac{\delta^2}{24\sqrt2}\|\mathsf A v\|_M.
\]
This bounds every subsequent bounded-effect probability, and remains valid
after a common sequential instrument including its classical record.
Here $e(t)$ enters $\ddot q+\mathsf A q=v/\delta+e(t)$, equivalently
through the additional Hamiltonian term $-\sum_i m_i e_i(t)q_i$.
\end{proposition}
\begin{proof}
Diagonalize the positive oscillator matrix in mass-normalized coordinates.
A linear force changes only the coherent displacement, with a central phase.
The exact forced mean at the pulse endpoint in each mode is
\[
 q(\delta/2)=\frac{1-\cos(\delta\omega)}{\delta\omega^2}v,
 \qquad \dot q(\delta/2)=\frac{\sin(\delta\omega)}{\delta\omega}v.
\]
These are the free evolution for time $\delta/2$ of position zero and
velocity $\operatorname{sinc}(\delta\omega/2)v$. Two such coherent states
with velocity difference $w$ at their common reference time have squared
overlap $\exp[-\|\mathsf A^{-1/4}w\|_M^2/2]$. Consequently their trace
distance is at most $\|\mathsf A^{-1/4}w\|_M/\sqrt2$.
The scalar inequality $|1-\operatorname{sinc}x|\le x^2/6$, together with
$\mathsf A\ge I$, gives the stated ideal error. For the additional force,
the displacement amplitude is a unitary mode rotation of
$\mathsf A^{-1/4}e(t)/\sqrt2$ integrated over the pulse. The triangle
inequality bounds its norm by $\epsilon_F/\sqrt2$. Unitary invariance
propagates the initial error, and contractivity under the common channel
gives the record-probability statement.
\end{proof}

The virtual reference state is not asserted to be the actual state halfway
through the force pulse. The comparison starts at the pulse endpoint.
The ideal energy added above the original vacuum is
$\|v_\delta\|_M^2/2\le\|v\|_M^2/2$. A trace-distance preparation bound
alone supplies no noisy-energy or apparatus-power bound.

On the same 64-site action, all sixteen nonzero entries of $v$ lie in the
original preparation region. Exact quadratic-field arithmetic gives
$\|v\|_M^2\in[0.552881768009,0.552881768010]$ and
$\|\mathsf A v\|_M^2\in[1965.046309667688,1965.046309667689]$.
For $10^{-4}\le\delta\le2\times10^{-3}$ and
$\epsilon_{\rm vac},\epsilon_F\le10^{-6}$, the conservative bound obtained
by dropping each factor $1/\sqrt2$ is less than $9.389\times10^{-6}$.
The baseline, with zero intended force, has error at most $2\times10^{-6}$.
Both fit the separate $10^{-5}$ preparation budgets of the finite-duration
instrument. Preparation ends more than $0.05148$ model-time units before
its first pointer operation. All 21 correlated readouts and all fifteen
resolved response signs therefore retain their stated bounds.

The operation contract includes 16 force ports, 48 coefficient reads and
32 on/off writes, together with the integrated force-error limit. It is a
specified quantum-control channel on the supplied action, not an
authenticated quantum execution or a derivation of its control law.
The independent verifier reconstructs positivity and mass self-adjointness
of all 64 action rows, the preparation norm and the inherited timing/error
interfaces. Altered pulse centers, omitted resets, incorrect coefficients,
missing branch errors and promoted outcome claims are rejected.
